\documentclass[9pt,twoside]{extarticle}
\usepackage{Breviary-Appendix}
\usepackage{indentfirst}
%%
\begin{document}
\renewcommand{\thepage}{\arabic{page}*}
\begin{center} \includegraphics[width=1 \textwidth]{20251031_131444495} \end{center}
%\frontmatter
\subsection*{\center \normalfont \textsc{die 12 decembris}}}

\medskip
\section*{\center \cloisterblack \LARGE \textred{B. M}ariæ \textred{V}irg. de \textred{G}uadalupe}
\subsubsection*{\center \boldred{DUPLEX MAJUS I CLASSIS}}
    \fancyhead[CO]{\textsc{\textred{Beatæ Mariæ Virginis de Guadalupe}}}
    \fancyhead[CE]{\textsc{\textred{Die 12 Decembris}}}
\begin{multicols}{2}

{\small \textred{Omnia de Communi Festorum B. Mariæ Virginis [161], præter ea quæ hic habentur propria.}}

\subsubsection*{\center \normalfont \textsc{IN I VESPERIS}}

\smallskip
\indent {\small \textred{Antiphonæ de Laudibus, Psalmi ut in I Vesperis de Communi [161]; Capitulum, Responsorium et Hymnus item de Communi.}}

\textred{℣.}~Non fecit táliter omni natióni. \textred{℟.}~Et judícia sua non manifestávit eis, allelúja.

\textred{Ad Magnif. Ant.} Elégi et sanctificávi * locum istum, ut sit ibi nomen meum, et permáneant óculi mei et cor meum ibi cunctis diébus.

\subsubsection*{\center \textred{Oratio}}
\lettrine{D}eus, qui sub beatíssimæ Vírginis Maríæ singulári patrocínio constitútos perpétuis benefíciis nos cumulári voluísti~: præsta supplícibus tuis; ut, cujus hódie commemoratióne lætámur in terris, ejus conspéctu perfruámur in cælis. Per Dóminum.

\small \textred{Fit Commemoratio Feriæ.}

\textred{Completorium de Dominica.}

\textred{In fine Hymnorum Completorii et Horarum dicitur:}\\
\textred{G}lória tibi, Dómine,\\
Qui natus es de Vírgine,\\
Cum Patre, et Sancto Spíritu,\\
In sempitérna sǽcula.

\textred{Ad Nunc dimittis Ant.} Allelúja, Sancta.

\normalsize

\subsubsection*{\center \normalfont \textsc{AD MATUTINUM}}

\smallskip
\indent \textred{Invitat.} In honórem beatíssimæ Maríæ Vírginis \textred{*} Jubilémus Dómino. \textred{Ps. 94} Veníte, exsultémus. \textred{2.}

\subsubsection*{\center \textred{Hymnus.}}
\lettrine{P}ræclára custos Vírginum,\\
Intácta Mater Núminis,\\
Cæléstis aulæ jánua,\\
Spes nostra, cæli gáudium.

\smallskip
\noindent \textred{I}nter rubéta lílium,\\
Colúmba formosíssima,\\
Virga e radíce gérminans\\
Nostro medélam vúlneri,

\smallskip
\noindent \textred{T}urris dracóni impérvia,\\
Amíca stella náufragis,\\
Tuére nos a fráudibus,\\
Tuáque luce dírige.

\smallskip
\noindent \textred{E}rróris umbras díscute,\\
Syrtes dolósas ámove,\\
Fluctus tot inter, déviis\\
Tutam reclúde sémitam.

\smallskip
\noindent \textred{G}lória tibi, Dómine,\\
Qui natus es de Vírgine,\\
Cum Patre, et Sancto Spíritu,\\
In sempitérna sǽcula.\\
{\null \hfill \textred{A}men.}

\subsubsection*{\center \normalfont \textsc{IN I NOCTURNO}}

\smallskip
\indent \textred{Ant.} A solis ortu * usque ad occásum, magnum est nomen meum in géntibus.

{\small \textred{Psalmi de Communi [165].}}

\textred{Ant.} Ego in altíssimis * habitávi, et thronus meus in colúmna nubis.

\textred{Ant.} Invénta sum * a non quæréntibus me; palam appárui iis, qui me non interrogábant.

\smallskip
\textred{℣.}~Notum fecit Dóminus salutáre suum. \textred{℟.}~In conspéctu géntium revelávit glóriam suam, allelúja.

\medskip
{\small \textred{Bened.} Alma Virgo vírginum intercédat pro nobis ad Dóminum.}
\subsubsection*{\center \normalfont De libro Ecclesiástici}
\subsubsection*{\textred{Lectio j \hfill Cap.~24,~5-31}}
\lettrineThree{E}{go} ex ore Altíssimi prodívi, primogénita ante omnem creatúram:
ego feci in cælis ut orirétur lumen indefíciens, et sicut nébula texi omnem terram:
ego in altíssimis habitávi, et thronus meus in colúmna nubis.
Gyrum cæli circuívi sola, et profúndum abýssi penetrávi, in flúctibus maris ambulávi,
et in omni terra steti~: et in omni pópulo,
et in omni gente primátum hábui:
et ómnium excelléntium et humílium corda virtúte calcávi~: et in his ómnibus réquiem quæsívi, et in hereditáte Dómini morábor.
Tunc præcépit, et dixit mihi Creátor ómnium~: et qui creávit me, requiévit in tabernáculo meo,
et dixit mihi~: In Jacob inhábita, et in Israël hereditáre, et in eléctis meis mitte radíces.

\textred{℟.}~Vidi speciósam sicut colúmbam ascendéntem désuper rivos aquárum~: cujus inæstimábilis odor erat nimis in vestiméntis ejus~: \textred{*} Et sicut dies verni circúmdabant eam flores rosárum, et lília convállium. \textred{℣.}~Quæ est ista, quæ ascéndit per desértum sicut vírgula fumi ex aromátibus myrrhæ et thuris? \textred{E}t sicut dies verni.

\medskip
{\small \textred{Bened.} Nos cum prole pia benedícat Virgo María.}
\subsubsection*{\center \textred{Lectio ij}}
\lettrine{A}{b} inítio, et ante sǽcula creáta sum, et usque ad futúrum sǽculum non désinam, et in habitatióne sancta coram ipso ministrávi.
Et sic in Sion firmáta sum, et in civitáte sanctificáta simíliter requiévi, et in Jerúsalem potéstas mea.
Et radicávi in pópulo honorificáto, et in parte Dei mei heréditas illíus, et in plenitúdine sanctórum deténtio mea.
Quasi cedrus exaltáta sum in Líbano, et quasi cypréssus in monte Sion:
quasi palma exaltáta sum in Cades, et quasi plantátio rosæ in Jéricho:
quasi olíva speciósa in campis, et quasi plátanus exaltáta sum juxta aquam in platéis.
Sicut cinnamómum, et bálsamum aromatízans odórem dedi~: quasi myrrha elécta dedi suavitátem odóris:
et quasi storax, et gálbanus, et úngula, et gutta, et quasi Líbanus non incísus vaporávi habitatiónem meam, et quasi bálsamum non mixtum odor meus.

\textred{℟.}~Quæ est ista, quæ ascéndit sicut auróra consúrgens, pulchra ut luna, elécta ut sol, \textred{*} Terríbilis ut castrórum ácies ordináta? 
\textred{℣.}~Fília Sion, tota formósa et suávis es, pulchra ut luna, elécta ut sol. \textred{T}erríbilis.

\medskip
{\small \textred{Bened.} Sancta Dei Génitrix sit nobis auxiliátrix.}
\subsubsection*{\center \textred{Lectio iij}}
\lettrine{E}{go} quasi terebínthus exténdi ramos meos, et rami mei honóris et grátiæ.
Ego quasi vitis fructificávi suavitátem odóris~: et flores mei fructus honóris et honestátis.
Ego mater pulchræ dilectiónis, et timóris, et agnitiónis, et sanctæ spei.
In me grátia omnis viæ et veritátis, in me omnis spes vitæ et virtútis.
Transíte ad me omnes qui concupíscitis me, et a generatiónibus meis implémini:
spíritus enim meus super mel dulcis, et heréditas mea super mel et favum:
memória mea in generatiónes sæculórum.
Qui edunt me, adhuc esúrient~: et qui bibunt me, adhuc sítient.
Qui audit me, non confundétur~: et qui operántur in me, non peccábunt.
Qui elúcidant me, vitam ætérnam habébunt.

\textred{℟.}~Quæ est ista, quæ procéssit sicut sol, et formósa tamquam Jerúsalem? \textred{*} Vidérunt eam fíliæ Sion, et beátam dixérunt, et regínæ laudavérunt eam. \textred{℣.}~Et sicut dies verni circúmdabant eam flores rosárum, et lília convállium. \textred{V}idérunt. Glória Patri. \textred{V}idérunt.

\subsubsection*{\center \normalfont \textsc{IN II NOCTURNO}}

\smallskip
\textred{Ant.} Signum magnum * appáruit in cælo~: múlier amícta sole, et luna sub pédibus ejus.

\textred{Ant.} Pópulus qui ambulábat * in ténebris, vidit lucem magnam~: habitántibus in regióne umbræ mortis, lux orta est eis.

\textred{Ant.} Vidérunt * omnes términi terræ salutáre Dei nostri.

\smallskip
\textred{℣.}~Introíbimus in tabernáculum ejus. \textred{℟.}~Adorábimus in loco ubi stetérunt pedes ejus, allelúja.

\medskip
\small \textred{Bened.} Précibus suæ Matris benedícat nos Fílius Dei Patris.

\normalsize
\subsubsection*{\center \textred{Lectio iv}}
\lettrineThree{A}{n}no a reparáta salúte millésimo quingentésimo trigésimo primo, Deípara Virgo Joánni Dídaco, pio rudíque neóphyto, Méxici in colle Tepejacénsi, uti antíqua et constánti traditióne mandátur, sese vidéndam prǽbuit, eúmque peramánter álloquens, epíscopum adíre, ipsíque ædem ínibi sibi sacram excitándam nuntiáre jussit. Rei veritátem solérter exploratúrus, responsiónem dístulit Joánnes de Zumárraga, loci antístes~: ac neóphyto, beatíssimæ Vírginis aspéctu imperióque rursus permóto, legationémque lácrimis et précibus iteránti, præcépit, ut signum, quo magnæ Dei Matris ostenderétur volúntas, eníxe péteret.

\smallskip
\textred{℟.}~Signum magnum appáruit in cælo, múlier amícta sole et luna sub pédibus ejus, \textred{*} Et in cápite ejus coróna duódecim stellárum.
\textred{℣.}~Fundaménta ejus in móntibus sanctis, díligit Dóminus portas Sion super ómnia tabernácula Jacob. \textred{E}t.

\medskip
\small \textred{Bened.} Jesus Maríæ Fílius sit nobis clemens et propítius.

\normalsize
\subsubsection*{\center \textred{Lectio v}}
\lettrine{D}{um} viam a colle Tepejacénsi remotiórem neóphytus arríperet, et Méxicum, ne pátruo letháli morbo corrépto extréma deéssent sacraménta, sacerdótem advocatúrus conténderet, eídem benigníssima Virgo tértio occúrrit; mæréntem de pátrui sanitáte certiórem facit; rosásque pulchérrimas, contra loci hiemísque asperitátem, recens obórtas, in ejus pállio compósitas, epíscopo deférri injúngit. Mandátis obséquitur Dídacus, cujus in pállio, rosis coram epíscopo effúsis, Maríæ sanctæ imágo qua prorsus spécie in colle prope urbem sese exhibúerat, mirum in modum depícta conspícitur. Tanto percúlsi prodígio religiósam cives effígiem in episcopáli sacéllo rite asservándam curant~: quæ, paulo post, solémni pompa ad ædem in collo Tepejacénsi exstrúctam transláta, insígni géntium cunctárum veneratióne præfúlget.

\smallskip
\textred{℟.}~Quæ est ista, quæ progréditur quasi auróra consúrgens~: \textred{*} Pulchra ut luna, elécta ut sol?
\textred{℣.}~Quasi arcus refúlgens inter nébulas glóriæ, et quasi flos rosárum in diébus vernis. \textred{P}ulchra.

\medskip
\small \textred{Bened.} Fílius Vírginis Maríæ det nobis gáudia vitæ.

\normalsize
\subsubsection*{\center \textred{Lectio vj}}
\lettrine{M}{a}gnífico deínceps excépta templo, cui canonicórum collégium Románi Pontífices ad divíni cultus splendórem addidérunt, Mexicánæ gentis pietátem in Deíparam summópere adáuxit; ingentíque abhinc quátuor sǽculis cólitur populórum ac miraculórum frequéntia. Quaprópter eam uti præsentíssimum advérsus públicas privatásque calamitátes præsídium archiepíscopus Mexicánus ceteríque illárum regiónum antístites, ómnium órdinum consensióne, in primáriam adlegérunt univérsæ Mexicánæ natiónis patrónam, ritéque eléctam Benedíctus décimus quartus auctoritáte apostólica declarávit, atque Offícium et Missam sub título Beátæ Vírginis Guadalupénsis recitári concéssit. Leo vero décimus tértius, iterátis Mexicanórum prǽsulum petitiónibus benígne ánnuens, novíssimum hoc Offícium, ex sacrórum Rítuum Congregatiónis consúlto, recitári indúlsit, Virginísque effígiem prodígiis atque cultu célebrem áurea coróna, suo nómine et jussu, solémni ritu condecorári decrévit. Quam quidem quarto Idus Octóbris anno millésimo octingentésimo nonagésimo quinto, coram epíscopis fere ómnibus reipúblicæ et áliis bene multis ex réliquis Américæ pártibus, máxima incolárum advenarúmque frequéntia et gratulatióne archiepíscopus Mexicánus pretiosíssimam impósuit.

\smallskip
\textred{℟.}~Elégi et sanctificávi locum istum~: \textred{*} Ut sit ibi nomen meum, et permáneant óculi mei, et cor meum ibi cunctis diébus.
\textred{℣.}~Non fecit táliter omni natióni, et judícia sua non manifestávit eis. \textred{U}t. Glória Patri. \textred{U}t.
\subsubsection*{\center \normalfont \textsc{IN III NOCTURNO}}

\smallskip
\textred{Ant.} Tabernáculum Dei * María, in médio civitátis ejus, et non commovébitur.

\textred{Ant.} Egréssa es * in salútem pópuli tui, in salútem cum Christo tuo.

\textred{Ant.} Gloriósa dicta sunt * de te, cívitas Dei~: fundávit te Dóminus in móntibus sanctis.

\smallskip
\textred{℣.}~Fílii tui de longe vénient. \textred{℟.}~Et fíliæ tuæ de látere surgent, allelúja.

\smallskip
\small \textred{Bened.} Fons Evangélii répleat nos dógmate cæli.

\normalsize
\subsubsection*{\center  \normalfont Léctio sancti Evangélii\\ secúndum Lucam}
\subsubsection*{\center \textred{Lectio vij \hfill Cap.~1,~39-47}}
\lettrine{I}{n} illo témpore~: Exsúrgens María ábiit in montána cum festinatióne, in civitátem Juda. Et intrávit in domum Zacharíæ, et salutávit Elísabeth. Et réliqua.

\subsubsection*{\center \normalfont Homilía sancti Bernárdi Abbátis\\
\small \textred{De verbis Apoc. cap. 12, ad finem}}

\smallskip
\lettrineThree{M}{i}rabátur Elísabeth persónam veniéntis, dicens~: Unde hoc mihi, ut véniat Mater Dómini mei ad me? Commendábat et vocem salutántis, adjíciens~: Ut facta est vox salutatiónis tuæ in áuribus meis, exsultávit in gáudio infans in útero meo~: et beatificábat fidem credéntis~: Beáta, ínquiens, quæ credidísti. Magna quidem præcónia; sed et devóta humílitas nihil sibi passa retinére, in eum magis univérsa refúdit, cujus in se benefícia laudabántur.

\textred{℟.}~Felix namque es, sacra Virgo María, et omni laude digníssima~: \textred{*} Quia ex te ortus est sol justítiæ, Christus, Deus noster.
\textred{℣.}~Ora pro pópulo, intérveni pro clero, intercéde pro devóto femíneo sexu~: séntiant omnes tuum juvámen, quicúmque célebrant tuam sanctam festivitátem. \textred{Q}uia.
 
\smallskip
\small \textred{Bened.} Ab hoste malígno deféndat nos Dei Génitrix Virgo.

\normalsize
\subsubsection*{\center \textred{Lectio viij}}
\lettrine{T}{u}, inquit, magníficas Matrem Dómini, sed magníficat ánima mea Dóminum. In voce mea fílium pérhibes exsultásse in gáudio, sed exsultávit spíritus meus in Deo, salutári meo~: et ipse quoque tamquam amícus sponsi gaudet ad vocem sponsi. Beátam me esse dicis, quæ crédidi; sed credulitátis et beatitúdinis causa respéctus est supérnæ pietátis; ut ex hoc magis beátam me dicant omnes generatiónes, quia ancíllam húmilem et exíguam respéxit Deus. Verúmtamen numquid putámus, fratres, Elísabeth sanctam in eo, quod per Spíritum útique loquebátur errásse? Absit. Beáta plane quam respéxit Deus, et beáta quæ crédidit. Hic enim magnus divínæ respectiónis éxstitit fructus.

\textred{℟.}~Beátam me dicent omnes generatiónes~: \textred{*} Quia fecit mihi Dóminus magna qui potens est, et sanctum nomen ejus.
\textred{℣.}~Et misericórdia ejus a progénie in progénies timéntibus eum. \textred{Q}uia.

\smallskip
\small \textred{Bened.} In omni tribulatióne et angústia succúrrat nobis pia Virgo María.

\smallskip
\textred{Si dicenda sit ix Lectio de Homilia Dominicæ, erit:}

\smallskip
\textred{Bened.} Verba Redemptóris pia sint suffrágia nobis.

\normalsize
\subsubsection*{\center \textred{Lectio ix}}
\lettrine{I}{n}effábili síquidem artifício Spíritus supervéniens tantæ humilitáti, magnanímitas tanta in secretário virgínei cordis accéssit, ut quod vidélicet nec humílitas tanta minúerit magnanimitátem, nec magnanímitas tanta humilitátem. Sed cum in sua æstimatióne tam húmilis esset, nihilóminus et in promissiónis credulitáte magnánimis, ut quæ nihil áliud quam exíguam sese reputábat ancíllam, ad inscrutábile sacraméntum nullátenus se dubitáret eléctam, et veram Dei et hóminis Genitrícem créderet mox futúram. Jam te, Mater misericórdiæ, per ipsum sinceríssimæ tuæ mentis afféctum tuis jacens provolúta pédibus luna, mediatrícem sibi apud solem justítiæ constitútam, devótis supplicatiónibus interpéllat, ut in lúmine tuo vídeat lumen, et solis grátiam tuo mereátur obténtu, quam vere amávit præ ómnibus, et ornávit stola glóriæ índuens, et corónam pulchritúdinis ponens in cápite tuo.

\textred{℟.}~Christi Virgo dilectíssima, virtútum operátrix, opem fer míseris~: \textred{*} Súbveni, Dómina, clamántibus ad te júgiter. \textred{℣.}~Quóniam peccatórum mole prémimur, et non est qui ádjuvet. \textred{S}úbveni. Glória Patri. \textred{S}úbveni.

Te Deum laudámus. \textred{6.}

\textred{Sacerd. ℣.}~Ora pro nobis, sancta Dei Génitrix. \textred{℟.}~Ut digni efficiámur promissiónibus Christi, allelúja.

\subsubsection*{\center \normalfont \textsc{AD LAUDES}\\ \normalfont \textred{et per Horas}}

\smallskip
\textred{Ant. 1} Quæ est ista * quæ procéssit sicut sol, et formósa tamquam Jerúsalem.

{\small \textred{Psalmi de Dominica \textsc{i} loco.}}

\textred{2} Vidérunt eam * fíliæ Sion vernántem in flóribus rosárum, et beatíssimam prædicavérunt.

\textred{3} Flores apparuérunt * in terra nostra~: te laudámus, sancta Dei Génitrix.

\textred{4} Cantáte ei * cánticum novum~: annuntiáte inter gentes glóriam ejus.

\textred{5} Tu glória * Jerúsalem, tu lætítia Israel, tu honorificéntia pópuli tui.

\subsubsection*{\center \normalfont \textred{Capitulum \hfill Eccli.~24,~23-24}}
\lettrine{E}{go} quasi vitis fructificávi suavitátem odóris~: et flores mei fructus honóris et honestátis. Ego mater pulchræ dilectiónis, et timóris, et agnitiónis, et sanctæ spei.

\subsubsection*{\center \normalfont \textred{Hymnus}}
\lettrine{A}{l}ma Parens, omni carens\\
Corruptélæ mácula,\\
Quam elégit qui confrégit\\
Diræ mortis víncula.

\smallskip
\noindent \textred{C}lemens, præsta, ut qui festa\\
Tua gaudent cólere,\\
Gratuléntur et læténtur\\
Veræ lucis lúmine.

\smallskip
\noindent \textred{P}rimum quidem nobis fidem\\
Tuis auge précibus,\\
Spe labéntes firma mentes,\\
Caritáte róbora.

\smallskip
\noindent \textred{M}ater Dei, nostræ spei\\
Causa et exórdium,\\
Aufer bella et flagélla,\\
Famem, pestem, gládium.

\smallskip
\noindent \textred{E}genórum et vinctórum\\
Consoláre gémitum~:\\
Fave votis, dans ægrótis\\
Optátum remédium.

\smallskip
\noindent \textred{P}ax et quies nostros dies\\
Fáciant lætíssimos~:\\
Inimícos fac amícos,\\
Seda malos hómines.

\smallskip
\noindent \textred{O} María, Mater pia,\\
Tuum da subsídium,\\
Ut regnémus et laudémus\\
Tuum semper Fílium.

\smallskip
\noindent \textred{P}ater Deus, Fili Deus,\\
Deus alme Spíritus,\\
Per ætérna nos gubérna\\
Sǽcula, Deus Trínitas.\\
{\null \hfill \textred{A}men.}

\smallskip
\textred{℣.}~Omnis terra adóret te, et psallat tibi. \textred{℟.}~Psalmum dicat nómini tuo, allelúja.

\smallskip 
\textred{Ad Bened. Ant.} Quasi arcus * refúlgens inter nébulas glóriæ, et quasi flos rosárum in diébus vernis.

\subsubsection*{\center \textred{Oratio}}
\lettrine{D}{e}us, qui sub beatíssimæ Vírginis Maríæ singulári patrocínio constitútos perpétuis benefíciis nos cumulári voluísti~: præsta supplícibus tuis; ut, cujus hódie commemoratióne lætámur in terris, ejus conspéctu perfruámur in cælis. Per Dóminum.

\smallskip
\small \textred{Et non fit Commemoratio Octavæ.}

\textred{Fit Commemoratio Feriæ.}

\textred{Ad Horas Antiphonæ de Laudibus, Psalmi de Dominica, ad Primam tamen ut in Festis et in Responsorio ℣.}~Qui de Vírgine dignátus es nasci. \textred{Reliqua ut in Communi [177].}

\normalsize

\subsubsection*{\center \normalfont \textsc{IN II VESPERIS}}

\smallskip
\small \textred{Antiphonæ de Laudibus, Psalmi, Capitulum et Hymnus} Ave, maris stella, \textred{de Communi [177].}

\normalsize \textred{℣.}~Non fecit táliter omni natióni. \textred{℟.}~Et judícia sua non manifestávit eis, allelúja.

\textred{Ad Magnif. Ant.} Leva in circúitu * óculos tuos et vide~: omnes isti congregáti sunt, venérunt tibi~: fílii tui de longe vénient, et fíliæ tuæ de látere surgent.

{\small \textred{Fit Commemoratio sequentis:}}

\textred{Ant.} In patiéntia tua * possedísti ánimam tuam, Lúcia, sponsa Christi~: odísti quæ in mundo sunt, et corúscas cum Angelis~: sánguine próprio inimícum vicísti.

\textred{℣.}~Diffúsa est grátia in lábiis tuis. \textred{℟.}~Proptérea benedíxit te Deus in ætérnum, allelúja.

\noindent Orémus. \hfill \textred{Oratio.}

\lettrine{E}{x}áudi nos, Deus, salutáris noster~: ut, sicut de beátæ Lúciæ Vírginis et Mártyris tuæ festivitáte gaudémus; ita piæ devotiónis erudiámur afféctu. Per Christum.

\smallskip
\small \textred{Et fit Commemoratio Feriæ.}

\textred{Completorium de Dominica.} 

\textred{Ad Nunc dimittis Ant.} Allelúja, Sancta.

\normalsize


\begin{center}\textcolor{red}{\rule{\linewidth}{0.5pt}} \end{center}
%\newpage
%\vspace{0.4\paperheight}

\subsubsection*{\center \normalfont IMPRIMATUR.\\
Tornaci, die 11 Februarii 1939.\\
\textsc{J. Lecouvet}, Vic. Gen.}

%\fancyfoot[C]{{\small Imprimé en Belgique par la Société S. Jean l'Evangéliste,
%Desclée & Cie, Tournai. - 5767}}
\fancyhead[C]{}
\fancyfoot[C]{}
\begin{center}\textcolor{red}{\rule{\linewidth}{0.5pt}} \end{center}
\end{multicols}
\end{document}